\documentclass[11pt]{article}
\usepackage[utf8]{inputenc}
\usepackage{amsmath,amssymb}
\usepackage{booktabs}
\usepackage{authblk}
\usepackage[margin=1in]{geometry}
\usepackage[hidelinks]{hyperref}
\usepackage{setspace}

\newcommand{\grad}{\nabla}
\newcommand{\half}{\tfrac12}

\title{The Hydrodynamic Form of RealQM:\\
\large Madelung Hydrostatics, a Surface Tension, and the Missing Half}

\author[1]{Claes Johnson}
\affil[1]{KTH Royal Institute of Technology, SE-100 44 Stockholm, Sweden \\
\texttt{claesjohnson@gmail.com}}
\date{\today \\[0.4em]
\small\textit{Mathematical formulation by the author; analysis developed} \\
\small\textit{with AI assistance (Claude, Anthropic).}}

\begin{document}
\maketitle

\begin{abstract}
\noindent RealQM models matter as non-overlapping unit charge densities in ordinary three-dimensional
space, one per electron, meeting at free boundaries and minimising the Coulomb energy. This note
records an exact structural fact about that construction which does not appear to have been stated:
\textbf{its kinetic term is the von Weizs\"acker functional}, and the functional derivative of the
von Weizs\"acker functional is \textbf{Bohm's quantum potential}. RealQM therefore already contains
the quantum potential, not as an analogy but as the variation of a term it writes down for other
reasons. Three consequences follow. First, the static theory is identified: minimising subject to
$\int\psi_i^2=1$ gives $Q+V=\mu$, quantum potential plus classical potential constant, which is the
\emph{hydrostatic equilibrium of a Madelung fluid} --- one fluid per electron, with free boundaries
where they meet. Second, the interface parameter that the electron-gas analysis leaves
underdetermined is identified physically: a Robin condition $\partial_n\psi+\beta\psi=0$ is the
natural boundary condition of a surface energy $\tfrac{\beta}{2}\oint\psi^2\,dS$, so $\beta$ is a
\textbf{surface tension of the electron--electron interface}, and the defect $C=0$ is the statement
that this surface tension is zero. That single identification explains why one parameter governs both
of the framework's outstanding difficulties: a surface tension simultaneously supplies confinement
energy and resists boundary motion, so it cannot be raised to obtain degeneracy pressure without
introducing resistance to transport. Third, the missing dynamics is named. What RealQM lacks is the
other half of the Madelung pair, continuity together with an Euler equation driven by
$-\grad(V+Q)$; supplying it per domain would furnish propagation, a velocity field carrying a
wavelength, and a boundary that travels, which are the three things separately blocked in the
conduction and interference problems. We also state what this is \emph{not}. It is not Bohmian
mechanics: there is no wavefunction on $3N$-dimensional configuration space and no point particle
besides the density, and configuration space is exactly where Bohm locates Bell-violating
correlations, so the entanglement question is not answered here and is not answered by adding
Madelung dynamics. It is closer to Schr\"odinger's own charge-density interpretation of 1926, from
which it differs by the free boundary and the rule of one domain per electron --- the two features
that meet the objections which caused Schr\"odinger to abandon it.
\end{abstract}

\setstretch{1.15}

\section{The identity}

RealQM \cite{realqm,realqmbook} represents an $N$-electron system by $N$ non-negative densities
$\psi_i$, each carrying unit charge on its own domain $\Omega_i$, the domains non-overlapping and
their interfaces free to move to where the energy is least. The energy is Coulomb's,
\begin{equation}\label{eq:E}
E[\{\psi_i\},\{\Omega_i\}]=\sum_i\int_{\Omega_i}\Bigl(\half|\grad\psi_i|^2
-\sum_a\frac{Z_a}{|x-x_a|}\psi_i^2\Bigr)dx
+\half\sum_{i\neq j}\iint\frac{\psi_i^2(x)\psi_j^2(y)}{|x-y|}\,dx\,dy,
\end{equation}
subject to $\int_{\Omega_i}\psi_i^2=1$, and the first term is the one of interest here. Write the
charge density of the $i$th electron as $\rho_i=\psi_i^2$. Then $\grad\rho_i=2\psi_i\grad\psi_i$, so
$|\grad\rho_i|^2/\rho_i=4|\grad\psi_i|^2$ and
\begin{equation}\label{eq:vW}
\half\int|\grad\psi_i|^2\,dx \;=\; \frac18\int\frac{|\grad\rho_i|^2}{\rho_i}\,dx
\;=\; T_{\rm vW}[\rho_i],
\end{equation}
in units $\hbar=m=1$. The right-hand side is the \textbf{von Weizs\"acker functional}. The
identification is elementary and exact; no approximation has been made and no limit taken.

Its consequence is not elementary. The functional derivative of $T_{\rm vW}$ is, with
$u=\sqrt{\rho}$,
\begin{equation}\label{eq:Q}
\frac{\delta T_{\rm vW}}{\delta\rho}
=\frac{1}{2u}\Bigl(-\grad^2u\Bigr)
=-\half\frac{\grad^2\sqrt{\rho}}{\sqrt{\rho}}
\;=\;Q,
\end{equation}
which is \textbf{Bohm's quantum potential}. RealQM therefore carries the quantum potential already,
not by resemblance but as the variation of a term it writes down for an entirely different reason ---
as the localisation cost of confining a charge to a region.

\subsection*{The two steps, in full}

\paragraph{Step one, \eqref{eq:vW}.} With $\rho=\psi^2$ and $\psi\ge0$, so that $\psi=\sqrt\rho$, the
chain rule gives $\grad\rho=2\psi\grad\psi$, hence
\[
|\grad\rho|^2=4\psi^2|\grad\psi|^2,\qquad
\frac{|\grad\rho|^2}{\rho}=\frac{4\psi^2|\grad\psi|^2}{\psi^2}=4|\grad\psi|^2,\qquad
\frac18\frac{|\grad\rho|^2}{\rho}=\half|\grad\psi|^2 ,
\]
and integrating gives \eqref{eq:vW}. There is no approximation and no limit in this; it is the chain
rule.

\paragraph{Step two, \eqref{eq:Q}.} Write $u=\sqrt\rho$, so $T_{\rm vW}=\half\int|\grad u|^2$, and
perturb $\rho\to\rho+\varepsilon\eta$. Since
$\sqrt{\rho+\varepsilon\eta}=u\bigl(1+\varepsilon\eta/\rho\bigr)^{1/2}
=u+\varepsilon\eta/(2u)+O(\varepsilon^2)$, the induced variation of $u$ is
\begin{equation}\label{eq:du}
\delta u=\frac{\eta}{2u}.
\end{equation}
Then
\begin{equation}\label{eq:vary}
\delta T_{\rm vW}=\int\grad u\cdot\grad(\delta u)\,dx
=-\int(\grad^2u)\,\delta u\,dx+\oint_{\partial\Omega}(\partial_nu)\,\delta u\,dS ,
\end{equation}
by parts, and substituting \eqref{eq:du},
\[
\delta T_{\rm vW}=\int\Bigl[-\frac{\grad^2u}{2u}\Bigr]\eta\,dx
+\oint_{\partial\Omega}(\partial_nu)\frac{\eta}{2u}\,dS .
\]
The coefficient of $\eta$ in the interior is the functional derivative, giving \eqref{eq:Q}.

\paragraph{The boundary term is the other half of the story.} The surface integral left over in
\eqref{eq:vary} is not a technicality to be discarded --- it is what fixes the interface condition,
and it is \S\ref{sec:surface} below. Requiring it to vanish for arbitrary $\eta$ gives
$\partial_nu=0$, the Neumann condition the framework presently uses. Adding a surface energy
$\tfrac{\beta}{2}\oint u^2\,dS$ contributes $\beta\oint u\,\delta u\,dS$, so the two boundary terms
combine to $\oint(\partial_nu+\beta u)\delta u\,dS$ and the natural condition becomes
$\partial_nu+\beta u=0$, Robin. \emph{The quantum potential and the surface tension emerge from the
same integration by parts}, one from the interior and one from the boundary.

\paragraph{A check by the other route.} Varying the $\tfrac18\int|\grad\rho|^2/\rho$ form directly,
\[
\delta T_{\rm vW}=\frac14\int\frac{\grad\rho\cdot\grad\eta}{\rho}
-\frac18\int\frac{|\grad\rho|^2}{\rho^2}\eta
=\int\Bigl[-\frac14\frac{\grad^2\rho}{\rho}+\frac18\frac{|\grad\rho|^2}{\rho^2}\Bigr]\eta\,dx,
\]
using $\grad\!\cdot(\grad\rho/\rho)=\grad^2\rho/\rho-|\grad\rho|^2/\rho^2$. Substituting $\rho=u^2$,
for which $\grad^2\rho=2|\grad u|^2+2u\grad^2u$ and $|\grad\rho|^2=4u^2|\grad u|^2$, the
$|\grad u|^2/u^2$ terms cancel exactly and what remains is $-\tfrac12\grad^2u/u$. The same answer by
a different path.

\section{What the static theory is}

Minimising \eqref{eq:E} in $\psi_i$ subject to $\int\psi_i^2=1$ gives the Euler--Lagrange equation
\begin{equation}\label{eq:EL}
-\half\grad^2\psi_i+V_i\psi_i=\mu_i\psi_i,
\end{equation}
where $V_i$ collects the nuclear attraction and the Hartree potential of the other domains, and
$\mu_i$ is the multiplier. Divide by $\psi_i$ and use \eqref{eq:Q}:
\begin{equation}\label{eq:hydrostatic}
\boxed{\;Q_i+V_i=\mu_i\;}
\end{equation}
Quantum potential plus classical potential constant on each domain. That is the condition of
\textbf{hydrostatic equilibrium} for a fluid whose internal energy is $T_{\rm vW}$: the chemical
potential is uniform, so nothing flows.

So RealQM in its present form is \emph{Madelung hydrostatics} --- one Madelung fluid per electron, in
ordinary three-dimensional space, the fluids meeting at free boundaries. It is a statics, which is
why it returns equilibrium partitions and no currents, and why the conduction study found that
``electrons have no equation of motion at all'' in it. The diagnosis is now sharper than that
sentence: it is not that a dynamics is missing in some unspecified way, but that exactly one half of
a known pair is present.

\section{The interface is a surface tension}\label{sec:surface}

The same view settles what the free interface is, and thereby what the framework's one
underdetermined number means.

Take a single domain and add to its energy a term proportional to the boundary integral of the
density,
\begin{equation}\label{eq:surf}
E_\sigma=\half\int_\Omega|\grad\psi|^2\,dx+\frac{\beta}{2}\oint_{\partial\Omega}\psi^2\,dS .
\end{equation}
Varying gives $\delta E_\sigma=\int(-\grad^2\psi)\delta\psi
+\oint(\partial_n\psi)\delta\psi+\beta\oint\psi\,\delta\psi$, so the natural boundary condition is
\begin{equation}\label{eq:robin}
\partial_n\psi+\beta\psi=0,
\end{equation}
the Robin condition exactly. Hence:

\begin{center}
\small
\begin{tabular}{@{}lll@{}}
\toprule
boundary condition & surface energy & fluid reading\\
\midrule
Neumann, $\beta=0$ (the present choice) & none & free surface, no surface tension\\
Robin, $\beta$ finite & $\tfrac{\beta}{2}\oint\psi^2dS$ & finite surface tension\\
Dirichlet, $\beta\to\infty$ & infinite & the density is pinned to zero at the face\\
\bottomrule
\end{tabular}
\end{center}

\noindent \textbf{$\beta$ is a surface tension of the electron--electron interface}, and the
framework has been running at zero surface tension. The electron-gas analysis, where the
localisation coefficient in $t(n)=C\,n^{5/3}$ comes out as $C=0$ under Neumann and
$C=C_{\rm TF}=2.871$ under Robin at $\beta a=1.146$, is then a measurement of that surface tension
against the degenerate electron gas.

\paragraph{Why one number governs two difficulties.} This identification explains something that was
otherwise a coincidence. A surface tension does two things at once, and they are the two things the
framework respectively lacks and enjoys:

\begin{center}
\small
\begin{tabular}{@{}ll@{}}
\toprule
it costs energy to \emph{have} interface area & confinement energy, hence degeneracy pressure\\
it costs energy to \emph{move} the interface & resistance to transport\\
\bottomrule
\end{tabular}
\end{center}

\noindent At $\beta=0$ the first is absent, which removes degeneracy pressure, screening and
bandwidth; and the second is absent, which is precisely why charge transport by boundary motion costs
nothing. These are not two independent defects to be repaired separately. They are one physical
quantity, read on its two faces, and no value of $\beta$ can supply the first without incurring the
second. Whether some intermediate value serves both is a quantitative question, and it is the
sharpest form of the open problem.

\section{The missing half}

The other half of the Madelung pair is the dynamics. For a Madelung fluid of density $\rho$ and
velocity $v$,
\begin{align}
\frac{\partial\rho}{\partial t}+\grad\!\cdot(\rho v)&=0,\label{eq:cont}\\
\frac{\partial v}{\partial t}+(v\cdot\grad)v&=-\grad\bigl(V+Q\bigr),\label{eq:euler}
\end{align}
with $Q$ from \eqref{eq:Q}. Setting $v=0$ in \eqref{eq:euler} returns \eqref{eq:hydrostatic}, so the
present theory is the equilibrium of this pair and nothing has to be given up to adopt it.

Supplying \eqref{eq:cont}--\eqref{eq:euler} per domain, with the interfaces evolving as free surfaces
between two such fluids --- the kinematic condition that the interface moves with the normal velocity,
and continuity of the chemical potential across it --- would furnish, in one step, the three things
found separately missing elsewhere in this programme:

\begin{itemize}
\item \textbf{Propagation.} A density that moves rather than relaxes, which is what a scattering or
interference problem requires and what an energy minimisation cannot supply.
\item \textbf{A wavelength.} The velocity field carries a phase, and with it the de Broglie relation;
a minimisation has no oscillation in it, its ground state having no nodes, so no repetition of static
calculations under perturbation can produce a fringe spacing.
\item \textbf{A boundary that travels.} Conduction in this framework is boundary motion, and
\eqref{eq:cont} is the equation that moves it.
\end{itemize}

\paragraph{What this costs, honestly.} Two qualifications belong with the proposal. Madelung
hydrodynamics is equivalent to the Schr\"odinger equation only for irrotational $v$ with quantised
circulation, $\oint v\cdot dl=2\pi n$; the vorticity-free condition and the quantisation are extra
requirements, not consequences, and they are where the ``quantum'' re-enters a formulation that
otherwise looks like fluid mechanics. And the free-surface problem for two Madelung fluids in contact
is a genuine free-boundary problem, whose well-posedness is not established here.

\section{What this is not: Bohmian mechanics}

The quantum potential is Bohm's, but the theory is not. Bohmian mechanics requires two objects that
RealQM does not have: a wavefunction on $3N$-dimensional configuration space, and a point particle in
addition to $\psi$, guided by it. RealQM has $N$ separate densities in ordinary three-dimensional
space and nothing besides them; the density \emph{is} the electron, and there is no hidden variable
riding it.

The difference is not cosmetic, and the cost is specific. Configuration space is exactly where
Bohmian mechanics locates entanglement: the guiding field of two particles is a function on
$\mathbb{R}^6$, and its non-factorisability is what produces correlations violating Bell's
inequality. A construction committed to three dimensions has no such object, and the Madelung
dynamics proposed above does not supply one --- each electron has its own velocity field in its own
$\mathbb{R}^3$. \emph{The entanglement question is therefore not answered here, and is not answered
by this proposal.} It remains the standing cost of the framework's central design choice, alongside
the Fermi surface.

\section{What this is not: Schr\"odinger 1926}

The nearest relative is not Bohm but Schr\"odinger's own interpretation of $\psi$ as a literal charge
density, proposed in 1926 and abandoned by him. It was abandoned for two reasons, and the framework's
two additional ingredients bear directly on them.

The first was \emph{spreading}: a free wave packet disperses without limit, so an electron would not
remain an electron. RealQM's answer is the free boundary together with non-overlap --- the domain is
bounded by its neighbours, and a partition of space does not disperse.

The second was \emph{localisation on measurement}: a spread density should deposit spread charge,
whereas an electron arrives at a screen as a single flash. Here the framework has a mechanism that
linear quantum mechanics does not, and it is worth stating precisely because it is the one place this
picture might do better rather than worse. Schr\"odinger's $\psi$ evolved linearly and could never
localise. RealQM's boundary moves so as to \emph{minimise the energy}, which is nonlinear and breaks
symmetry: a density arriving spread across a row of detector atoms is not at a minimum, and the
minimising partition is plausibly the one with the charge on a single atom, an attractive centre
being where the energy is. That would be a dynamical localisation rather than a postulated collapse.

This is a candidate mechanism and not a result, and the evidence available cuts both ways. In the
conduction study the behaviour that repeatedly disrupted the calculations was a domain seizing a
nuclear region outright, since with free interfaces no confinement cost opposes it --- which is the
wanted behaviour here, appearing as a pathology there. But the same study found the framework
\emph{over}-delocalises around neutral atoms, an added electron relaxing to a uniform corona at
$2.2\,a_0$ and still at $4.0\,a_0$. Whether an attractive site with a vacancy localises it is a
static calculation, entirely within what the solver already does, and it has not been performed.

\section{What would settle it}

Three computations, in increasing order of difficulty, and none of them requires new physics:

\begin{enumerate}
\item \textbf{The surface tension.} Test $\beta a=1.146$ on $\mathrm{H_2}$, where the domains meet at
the midplane \emph{between} the nuclei in a smooth region --- unlike helium, where the interface
passes through the nuclear cusp and interface error cannot be separated from cusp error. If binding
survives at the value the electron gas requires, the framework has a surface tension it has not yet
identified; if it does not, it has two incompatible ones.
\item \textbf{Localisation on a screen.} A spread density over a row of atoms, with and without an
attractive site carrying a vacancy. Does the minimising partition leave the charge spread, or collect
it on one atom? This decides whether the framework has a detection mechanism.
\item \textbf{The dynamics.} Implement \eqref{eq:cont}--\eqref{eq:euler} on a single domain and check
the free-particle case against a known wave packet, before attempting two domains with a moving free
surface between them.
\end{enumerate}

\noindent The first two are static and use machinery that exists. The third is the substantial one,
and it is what the rest of the programme is waiting on.

\begin{thebibliography}{9}
\bibitem{realqm} C.\ Johnson, \emph{RealQM: Quantum Mechanics as Multiphase 3D Continuum Mechanics},
in~\cite{realqmbook}.
\bibitem{realqmbook} C.\ Johnson, \emph{Real Quantum Mechanics} (Body\&Soul, Vol.\ VI),
\url{https://claes542.github.io/RealMolecule/realquantum3.pdf}.
\bibitem{gas} C.\ Johnson, \emph{Conduction in RealQM: the Free Boundary as Charge Carrier},
in~\cite{realqmbook}.
\end{thebibliography}

\end{document}
